% The big-M switch x_r <= M y_r and its LP relaxation, for two values of M.
%
%   figures/tikz/ip-bigm-relaxation.tex
%       -> media/figures/ip-bigm-relaxation.{png,svg}
%
% Lecture 3 (Integer Programming), section III "Reactor selection: the MILP",
% subsection "Logical constraints", read by the activity "What happens if we
% choose M = 20,000? Is the model still correct?" (panels (1)-(2)).
%
% The (x_r, y_r) plane with x_r >= 0, for one reactor r of eq:bigM.
%   Integer-feasible set: y_r = 0 forces x_r = 0 (the point (0,0)), and
%   y_r = 1 allows 0 <= x_r <= M (a segment on the line y_r = 1).
%   LP relaxation (y_r in [0,1]): the region 0 <= x_r <= M y_r, 0 <= y_r <= 1,
%   i.e. the triangle (0,0), (M,1), (0,1).
% SCALE: x_r in [0, 25] kmol/hr at 0.22 cm per unit (5.5 cm); y_r in [0, 1] at 2.4 cm.
% PANEL (1), M = 20: triangle vertices (0,0), (20,1), (0,1) -> (0,0), (4.4,2.4), (0,2.4) cm.
% PANEL (2), M = 20,000: the lower edge y_r = x_r/M rises to 25/20000 = 0.00125
%   at x_r = 25, i.e. 0.003 cm: drawn on the axis. Inside the window the relaxation
%   is the whole box except that sliver; the segment y_r = 1 runs off the window
%   (it ends at x_r = 20,000).
% The dashed line x_r = 12.5 is the flow of the optimal design. Where it meets the
%   lower edge (y_r = 0.625 in (1), 6.25e-4 in (2)) is the activity's answer, so the
%   figure does NOT mark or label that point.
% GREYSCALE: the relaxation is hatched, the integer-feasible set is thick black
% (dot and segment), the lower edge is a thin solid line and x_r = 12.5 is dashed.
% No distinction relies on colour; the one colour (blue hatch) is labelled.
\usetikzlibrary{patterns}
\definecolor{okabeblue}{HTML}{0072B2}

\begin{tikzpicture}[
    >={Latex[length=2.0mm]},
    font=\small,
    xs/.style={x=0.22cm, y=2.4cm},
    pt/.style={fill=black, draw=black, circle, inner sep=0pt, minimum size=4.2pt},
    lab/.style={font=\scriptsize, inner sep=1.5pt},
    panelcap/.style={font=\small, anchor=north},
  ]
% ---------------------------------------------------------------- panel (1)
\begin{scope}[xs]
  \path[pattern=north east lines, pattern color=okabeblue!60] (0,0) -- (20,1) -- (0,1) -- cycle;
  \draw[->, black!70] (0,0) -- (27,0) node[right, lab] {$x_r$};
  \draw[->, black!70] (0,0) -- (0,1.25) node[above, lab] {$y_r$};
  \draw (0,1) -- ++(-0.5,0) node[left, lab] {$1$};
  \node[left, lab] at (-0.3,0) {$0$};
  \draw (20,0) -- ++(0,-0.04) node[below, lab] {$M = 20$};
  \draw[thin] (0,0) -- (20,1);
  \node[lab, rotate=28.6, anchor=north] at (6.5,0.31) {$x_r = M y_r$};
  \draw[dashed, thick] (12.5,0) -- (12.5,1.08) node[above, lab] {$x_r = 12.5$};
  \draw[line width=2.4pt] (0,1) -- (20,1);
  \node[pt] at (0,0) {};
  \node[lab, fill=white, inner sep=1pt] at (5.2,0.80) {LP relaxation};
  \node[panelcap] at (12.5,-0.32) {(1) $M = 20$};
\end{scope}
% ---------------------------------------------------------------- panel (2)
\begin{scope}[xs, xshift=7.6cm]
  \path[pattern=north east lines, pattern color=okabeblue!60] (0,0) -- (25,0) -- (25,1) -- (0,1) -- cycle;
  \draw[->, black!70] (0,0) -- (27,0) node[right, lab] {$x_r$};
  \draw[->, black!70] (0,0) -- (0,1.25) node[above, lab] {$y_r$};
  \draw (0,1) -- ++(-0.5,0) node[left, lab] {$1$};
  \node[left, lab] at (-0.3,0) {$0$};
  \draw[dashed, thick] (12.5,0) -- (12.5,1.08) node[above, lab] {$x_r = 12.5$};
  \draw[line width=2.4pt, ->] (0,1) -- (26.5,1);
  \node[lab, anchor=south east] at (26.5,1.0) {to $M = 20{,}000$};
  \node[pt] at (0,0) {};
  \node[lab, fill=white, inner sep=1pt] at (5.2,0.55) {LP relaxation};
  \node[lab, anchor=north west, align=left] at (14.0,-0.04) {$x_r = M y_r$ lies\\[-2pt]on the axis};
  \node[panelcap] at (12.5,-0.32) {(2) $M = 20{,}000$};
\end{scope}
% ---------------------------------------------------------------- legend
\begin{scope}[yshift=-1.55cm]
  \draw[line width=2.4pt] (0.5,0) -- (1.3,0);
  \node[pt] at (1.6,0) {};
  \node[lab, anchor=west] at (1.8,0) {integer feasible ($y_r \in \{0,1\}$)};
  \path[pattern=north east lines, pattern color=okabeblue!60] (7.6,-0.12) rectangle (8.4,0.12);
  \draw[black!50] (7.6,-0.12) rectangle (8.4,0.12);
  \node[lab, anchor=west] at (8.5,0) {LP relaxation ($0 \le y_r \le 1$)};
\end{scope}
\end{tikzpicture}
